\section{Introduction}
Questions about robbery are often framed as when and where it happens. Police-reported data can answer a different and narrower question, and answering it carefully matters because two quite different goals share one dataset.
\emph{Inference} asks which features are associated with the outcome and by how much; \emph{prediction} asks how well the outcome can be forecast for new cases \citep{shmueli2010explain,breiman2001two}. A feature can be strongly and reliably associated with an outcome and still help little in predicting it, and the reverse.

We state the task in words that the rest of the paper keeps: \emph{given that a major property-type crime was reported, which type is it?} The outcome is whether a reported incident is a robbery rather than an auto theft, a break and enter or a theft over. This is $P(\text{robbery}\mid\text{reported incident})$. It is not the probability that a person is robbed at a place or time, which would need exposure data such as foot traffic or the number of people present, and it carries no causal reading. Reporting itself varies by crime type and neighbourhood \citep{skogan1977dark,lum2016predict}.

The paper builds on two earlier course projects on Toronto Police open data, a robbery classifier and a logistic-regression analysis of robbery. We re-ran both, corrected their problems and put them on common ground. Our contributions are:
\begin{itemize}
  \item a corrected binomial model for the odds of robbery, with hour, month and weekday as categorical terms, solar darkness, a post-2019 period term and neighbourhood-cluster-robust intervals;
  \item an ablation ladder of five models on identical rows, evaluated by neighbourhood-grouped cross-validation and by a temporal holdout, with bootstrap intervals for every metric and every rung-to-rung difference; and
  \item a direct contrast between the two analyses on the same data: many time terms are statistically clear, yet time alone adds little to prediction.
\end{itemize}
